\documentclass[10pt,a4paper]{amsart}
\usepackage[margin=19mm]{geometry}
\usepackage{amsmath,amssymb,mathtools}
\usepackage[scr=rsfs,cal=euler]{mathalfa}
\usepackage{xfrac}
\usepackage[unicode,hidelinks,bookmarks=false]{hyperref}
\newcommand{\ie}{i.e.\ }
\newcommand{\cf}{cf.\ }
\newcommand{\resp}{respectively\ }
\newcommand{\Subsection}{\subsection}
\providecommand{\nobreakdash}{\nobreak\hbox{-}}
\setlength{\emergencystretch}{2em}
\multlinegap0pt
\usepackage{bm}
\usepackage{bbm}
\usepackage{dsfont}
\usepackage{xspace}
\usepackage{cancel}
\usepackage{mathtools}
\usepackage{amsthm}
\makeatletter
\@ifundefined{theorem}{\newtheorem{theorem}{Theorem}}{}
\@ifundefined{lemma}{\newtheorem{lemma}[theorem]{Lemma}}{}
\makeatother
\newcommand{\Ec}{\mathbf E}
\newcommand{\Kc}{\mathbf K}
\newcommand{\M}{\operatorname{AG}}
\newcommand{\dd}{\,\mathrm d}
\title[Global convergence and monotonicity of Newton iteration for ]{Global convergence and monotonicity of Newton iteration for the inverse Gr\"otzsch modulus}
\author{Deguang Zhong}
\date{Source: arXiv:2608.22777v1 (CC BY 4.0)}
\begin{document}
\maketitle

\noindent\emph{Source and licence.} Global convergence and monotonicity of Newton iteration for the inverse Gr\"otzsch modulus. Version arXiv:2608.22777v1 (CC BY 4.0), distributed under the Creative Commons Attribution 4.0 International licence, as recorded in the arXiv licence metadata for this exact version and shown on the abstract page https://arxiv.org/abs/2608.22777v1. Full rights notice: licenses/arxiv-2608.22777-CC-BY-4.0.txt.\par\smallskip

\noindent\textit{Editorial scope.} Counted unit: the Lemma whose source label is \texttt{lem:inflection} in the article (source0.tex lines 178--193, source label \texttt{lem:inflection}), delivered together with the article's own complete original proof of it (source0.tex lines 195--274, one \texttt{proof} environment of 80 lines). No mathematics is written, repaired or re-derived: statement and proof are the article's own text line for line. Required context retained uncounted: eq:KE-derivatives (lines 126--130); eq:legendredcsdcs (lines 132--135); eq:mu-prime (lines 138--141). Every definition, display and label of those lines is retained unchanged. Omitted from this package and not redistributed: 1--125, 131--131, 136--137, 142--177, 194--194, 275--489. Nothing inside a retained excerpt is omitted or altered: every retained body is delivered exactly as the article's lines are written, so no omission receipt of any kind accompanies this package. Citations: the retained excerpts carry no citation command at all, so no bibliography entry of the article is transcribed and no reference is redistributed. Reference adaptation: none was needed and none was made; every reference inside the delivered unit resolves inside it, so no unresolved reference or citation is produced. Printed slips kept literal: no printed word, symbol, hypothesis or number of the counted statement or of its proof was altered. Layout and class: the article's own class is \texttt{article} with packages including fontenc, lmodern, amsmath, amssymb, amsthm, mathtools, geometry, microtype, hyperref; the wrapper is the lane's pinned \texttt{amsart} editorial wrapper at 10pt on A4, which restates the theorem-like environments the retained text needs and adds the article's own macro definitions that the retained text needs (\texttt{\textbackslash Ec}, \texttt{\textbackslash Kc}, \texttt{\textbackslash M}, \texttt{\textbackslash dd}), with extra wrapper packages (bm, bbm, dsfont, xspace, cancel, mathtools). The delivered numbering is the unit's own. Assets: no retained span contains a figure environment, a graphics-inclusion command, a PDF-page inclusion, a table or a TikZ picture; no image file is copied into this package, the package declares no asset dependency and no image file is redistributed with it. Licence: version arXiv:2608.22777v1 of this article is distributed under the Creative Commons Attribution 4.0 International licence, as recorded in the arXiv licence metadata for this exact version and shown on the abstract page https://arxiv.org/abs/2608.22777v1. A full rights notice is delivered at licenses/arxiv-2608.22777-CC-BY-4.0.txt. Characters: every retained excerpt is pure ASCII, so the delivered engine, XeLaTeX with its default Latin Modern text font, reports no missing character.
\begin{equation}\label{eq:KE-derivatives}
 \Kc'(r)=\frac{\Ec(r)-r'^2\Kc(r)}{rr'^2},
 \qquad
 \Ec'(r)=\frac{\Ec(r)-\Kc(r)}r,
\end{equation}

\begin{equation}\label{eq:legendredcsdcs}
 \Kc(r)\Ec(r')+\Ec(r)\Kc(r')
 -\Kc(r)\Kc(r')=\frac{\pi}{2}.
\end{equation}

\begin{equation}\label{eq:mu-prime}
 \mu'(r)=-\frac{\pi^2}{4rr'^2\Kc(r)^2}
 =-\frac{\M(1,r')^2}{rr'^2}<0,
\end{equation}

\begin{lemma}\label{lem:inflection}
There exists a unique $a\in(1/2,1/\sqrt2)$ such that
\begin{equation}\label{eq:convex-concave}
 \mu''(r)>0\quad(0<r<a),
 \qquad
 \mu''(r)<0\quad(a<r<1).
\end{equation}
Moreover, if $c=1/\sqrt2$ and
\[
 N_y(r)=r-\frac{\mu(r)-y}{\mu'(r)},
\]
then
\begin{equation}\label{eq:key-barrier}
 N_{\pi/2}(a)<1.
\end{equation}
\end{lemma}

\begin{proof}
Write
\[
 h(r)=rr'^2\Kc(r)^2,
 \qquad
 u(r)=\frac{\Ec(r)}{\Kc(r)},
 \qquad
 q(r)=2u(r)-(1+r^2).
\]
Equations \eqref{eq:KE-derivatives} give
\begin{equation}\label{eq:h-log}
 \frac{h'(r)}{h(r)}=\frac{q(r)}{rr'^2}.
\end{equation}
Since $\mu'=-\pi^2/(4h)$, it follows that
\begin{equation}\label{eq:musecond-sign}
 \operatorname{sgn}\mu''(r)=\operatorname{sgn}q(r).
\end{equation}
On the other hand, differentiating $u=\Ec/\Kc$ and using
\eqref{eq:KE-derivatives}, we find
\begin{equation}\label{eq:qprime}
 q'(r)=-\frac{2\{(u(r)-r'^2)^2+2r^2r'^2\}}{rr'^2}<0.
\end{equation}
Thus $q$ is strictly decreasing, and $q(0^+)=1$.

We next locate the zero of $q.$  Since $\Ec(r)-r'^2\Kc(r)>0$, one has $u(r)>r'^2$.  At $r=1/2$ this
gives
\[
 q(1/2)>2\cdot\frac34-\frac54=\frac14>0.
\]
At $c=1/\sqrt2$, the Legendre's relation (\ref{eq:legendredcsdcs}) reduces to
\[
 2\Kc(c)\Ec(c)-\Kc(c)^2=\frac\pi2,
\]
and hence
\[
 q(c)=\frac{\pi}{2\Kc(c)^2}-\frac12<0.
\]
For completeness, the last inequality follows already from the first three
positive terms of the hypergeometric series:
\[
 \Kc(c)>
 \frac\pi2\left(1+\frac18+\frac9{256}\right)
 =\frac{297\pi}{512}>\sqrt\pi.
\]
Therefore $q$ has a unique zero $a\in(1/2,c)$, and
\eqref{eq:convex-concave} follows from \eqref{eq:musecond-sign}.

It remains to prove the inequality \eqref{eq:key-barrier}.  Because
$\mu(c)=\pi/2$ and $\mu''<0$ on $(a,c)$, the function
$|\mu'|$ is increasing there.  Consequently,
\begin{equation}\label{eq:delta-bound}
 \mu(a)-\frac\pi2
 =\int_a^c|\mu'(t)|\dd t
 \le(c-a)|\mu'(c)|.
\end{equation}
Furthermore, the equation \eqref{eq:mu-prime} gives
\[
 \frac{|\mu'(c)|}{|\mu'(a)|}=\frac{h(a)}{h(c)}<2.
\]
Indeed, the elementary bounds
\[
 \frac\pi2<\Kc(r)<\frac\pi{2r'}
\]
imply
\[
 h(a)<\frac{\pi^2a}{4},
 \qquad
 h(c)>\frac{\pi^2c}{8},
 \qquad
 \frac{h(a)}{h(c)}<\frac{2a}{c}<2.
\]
Combining this with \eqref{eq:delta-bound}, we obtain
\[
 N_{\pi/2}(a)
 =a+\frac{\mu(a)-\pi/2}{|\mu'(a)|}
 <a+2(c-a)=2c-a
 <\sqrt2-\frac12<1,
\]
where we used $a>1/2$.  This proves \eqref{eq:key-barrier}.
\end{proof}

\end{document}
